\section{Self-testing is a typed extraction claim}
\label{sec:selftest}

Classical input/output data can constrain untrusted quantum apparatus up to
unavoidable freedoms. The original self-testing formulation compares devices
with a reference realization through local changes of representation
\cite{mayers2004}; modern definitions distinguish state and measurement
certification and quantify robustness \cite{supic2020}.

\subsection{State extraction and the right quantifiers}
Fix a resource class of tensor-product strategies, a game $G$, a target pure
state $\ket\phi$, and a threshold $u$. Define target overlap after local
extraction by
\begin{equation}
 F_{\rm ext}(\rho,\phi)=
 \sup_{\Lambda_1,\ldots,\Lambda_n}
 \bra\phi(\Lambda_1\otimes\cdots\otimes\Lambda_n)(\rho)\ket\phi,
 \label{eq:extraction}
\end{equation}
where each $\Lambda_i$ is a local completely positive trace-preserving map
into the corresponding target system. A robust state-extraction statement is
\begin{equation}
 \forall (\rho,M)\in\mathcal R:\quad
 S_G(p_{\rho,M})\geq u-\varepsilon
 \ \Longrightarrow\ F_{\rm ext}(\rho,\phi)\geq1-f(\varepsilon).
 \label{eq:robust}
\end{equation}
The channels can depend on the unknown realization; the target and function
$f$ cannot. An exact self-test has $u=\val_{\mathcal R}(G)$, $f(0)=0$,
and a nonvacuous optimal realization. A robust self-test requires
$f(\varepsilon)\to0$ as $\varepsilon\to0$. The fidelity convention in
Eq.~\eqref{eq:extraction} is the overlap, not its square root.

Local isometries give a representation-level description:
the physical state maps to a target tensor an unobserved auxiliary state.
Auxiliary entanglement, unused Hilbert-space sectors, and local basis choices
are permitted. Measurement certification additionally compares the action of
physical operators on the tested state with target operators under the same
local maps. Operator-norm equality on the entire physical Hilbert space is
generally unjustified because the experiment does not probe unused sectors.

Pairwise convertibility of two optimal physical realizations is therefore
not the definition of a unique self-tested target. Both may contain a common
extractable target with different auxiliary systems. To disprove extraction
of a proposed $\phi$, one must exhibit an admissible optimum for which
$F_{\rm ext}<1$, or establish another obstruction to the stated equivalence.
Distinct numerical state vectors, Schmidt spectra before removal of auxiliary
factors, or inequivalent measurement choices alone do not settle this question.

\begin{proposition}[A common target with incomparable realizations]
Let $\ket{\Phi_d}=d^{-1/2}\sum_{j=0}^{d-1}\ket{jj}$. The states
$\ket{\Phi_6}$ and $\ket{\Phi_{10}}$ both attain the optimal CHSH score, but
neither can be transformed into the other by a product of local trace-preserving
quantum channels. Both exactly extract $\ket{\Phi_2}$.
\end{proposition}
\begin{proof}
Relabeling a local basis factors $\Phi_6$ as $\Phi_2\otimes\Phi_3$ and
$\Phi_{10}$ as $\Phi_2\otimes\Phi_5$. Use the CHSH measurements on the first
factor and discard the second for extraction. If product local channels
transform a pure $\Phi_d$ into a pure $\Phi_m$, their local Stinespring
isometries transform it into $\Phi_m\otimes\eta$ for a pure environment state
$\eta$. Local isometries preserve Schmidt rank, hence $d=mr$ with
$r$ the Schmidt rank of $\eta$. Neither $6/10$ nor $10/6$ is an integer.
\end{proof}
This statement concerns product local channels without communication;
LOCC convertibility is a different relation.

\subsection{Complex conjugation and information loss}
For a standard classical-interface Bell experiment, simultaneous complex
conjugation of $\rho$ and all local POVM matrices preserves every probability:
the Born probability is real and conjugation conjugates its trace. Complex
conjugation is not in general a physical channel. Measurement targets with
complex entries must therefore specify the tolerated conjugation freedom;
flagged simulations can require more than a single unflagged choice
\cite{mckague2010}. A compiler must not promise unique complex apparatus
merely because a numerical optimizer chose one representative.

There are three distinct uniqueness questions. A scalar Bell score can expose
a unique behavior; a behavior can still admit many physical realizations;
and these realizations may or may not extract a common target under the
chosen equivalence. Strict convexity or uniqueness of an SDP moment matrix
does not establish the last property. A finite moment relaxation may also
contain points without a corresponding physical strategy.

\subsection{What can be automated responsibly}
An SOS certificate $u\id-\mathcal B=\sum_kP_k^\dagger P_k$ implies
$P_k\ket\psi=0$ at exact saturation. At score deficit $\varepsilon$,
\[
 \sum_k\|P_k\ket\psi\|^2\leq\varepsilon.
\]
These relations are useful inputs to a proof of an extraction map; they are
not themselves such a map. Analytic operator inequalities can bound extraction
fidelity directly for important games \cite{kaniewski2016}. A further route
is to choose an extraction ansatz and lower-bound its fidelity through moment
relaxations, but the ansatz, operator identities, and bound direction must all
be checked.

The initial toolchain verifies a CHSH SOS identity and evaluates reference
strategies. It does not claim to derive or verify a general self-testing
theorem. A future certificate format should record the target, allowed
dimensions, extraction maps or a proved construction, gauge conventions,
distance metric, score normalization, robustness function, and a proof
dependency list. This makes ``self-testing supported'' a precise capability
rather than a boolean attached to a game name.
